\subsection{Experience and Lessons Learned}
\label{emt:sec:lessons}

\noindent
Developing EMT-Linux and the MMU drivers together with 
    the emulated MMU
    took extensive engineering efforts.
Specifically, we had to develop interacting moving parts 
    across the hardware-software boundary.

A key principle
    is to enable incremental, contiguous engineering practice.
For example, when implementing ECPT, 
    our first milestone is to write a basic per-process hashed page table (BHPT)
    using EMT, without features like collision resolution,
    elastic resizing, CWT/CWC, etc.
We implement BHPT logic in both EMT-Linux and the MMU and have a running system 
    that only supports a tiny microbenchmark.
Despite being basic, the running system serves as a foundation for gradually
    adding features and eventually evolving into a full-fledged system.
For features that need both hardware and software support,
    we start from the hardware, which typically has simpler logic
    and mostly reads translation data.

We continuously test and evaluate our system, not only 
    for correctness (\S\ref{emt:ss:correctness}) but also for performance.
% For example, % when refactoring Linux to EMT-Linux (\S\ref{emt:sec:impl}),
Continuous profiling helped us quickly observe performance regression over vanilla Linux;
    when developing ECPT, profiling helped us identify major inefficiencies compared 
    to the radix-page-table-based system. 
We developed the instruction analysis 
    with flame graphs~\cite{tool:flamegraph} in our emulator toolchain. % (\S\ref{emt:sec:eval:emulator}) 
    and use them extensively. % during the development of the ECPT MMU driver.
We also added GDB support for the ECPT-based system
% \schai{ to accelerate debugging with its functionalities like breakpoint, tracing, and watchpoint. 
    (QEMU's GDB was hardcoded with radix).
% ; we changed it reuse
 %   the existing software mmu logic to acquire the physical address.}

We largely reuse the existing Linux compiler toolchain to keep the effort manageable.
We leverage x86's model-specific registers (MSR)~\cite{intel:man} as the control registers for ECPT,
    which is supported by GCC (x86's read/write MSR instruction allows access 
    to arbitrary MSR with a 32-bit identifier).
% \schaic{If we need a number of registers we added, it's 54.}
We use Clang's static analysis tools (e.g., clang-query~\cite{tool:clang:query}) to search for
    error-prone code patterns during refactoring.

% For architecture design like ECPT, the communication between hardware and kernel can be complicated:
    % 9 registers for kernel, 9 registers for user page table, and also associated register for CWT and rehashed way.

One mistake we made was to start from Linux's boot-time kernel page table~\cite{linux:head64s, linux:head64c},
    which is legacy code without debugging tools.
Developing boot-time kECPT requires implementation in assembly
    and build mappings 
    of kernel code/data
    from a bootstrap address space.
% However, dealing with them at the beginning can greatly slow down the development
  %  since the booting stage is full of legacy code 
% but running short of debugging tools.
% For example, perparing ECPT's kernel page table 
%    requires to prepare special 
% requires
%     prepare special mappings for kernel text and data;
%     it has to deal with x86's legacy booting process 
%         like operating at different address space before entering x86 long mode.
We spent four months to understand the details of the boot-time kernel page table.
Retrospectively, we should have started from runtime kECPTs
    (by switching the system to it after boot~\cite{linux:cpucommon})
    so we could have a running system quicker to parallelize our efforts.

% Even with the help of \branding, implementing a kernel support for a new translation architecture